A dual-process chess agent pairs a fast learned move scorer (System~1) with slow tree search (System~2). We ask whether System~1's own confidence can decide \emph{when} to search. Our agent uses LayaChess, a 421M-parameter model that scores every legal move, as System~1 and a PUCT search with a 2.3M-parameter value network as System~2. Move quality is judged by Stockfish~16 and playing strength by tournaments against strength-limited Stockfish anchors. In the pre-specified first stage (V1; \nPos{} positions, 500 games), always searching lowers average centipawn loss (ACPL) from \acplLaya{} to \acplSTwo{} (\dAcpl{}\cp, 95\% CI [\dAcplLo, \dAcplHi]) at \costRatio$\times$ the time per move, but the confidence gate does not beat a random gate at the same search rate (advantage \advFifty{}\cp\ [\advFiftyLo, \advFiftyHi] at \rateFifty\%). Two further predictions are reversed, the hybrid is not stronger than System~1 alone (Elo difference \eDHL{} [\eDHLLo, \eDHLHi]), and the value network returns a constant output on \constPct\% of test positions. Two follow-up stages on fresh test months changed the agent. V2 adds a material blunder filter, tablebases, a better search and a retrained value network: ACPL falls from 114.5 to 68.7 and Elo rises by about 200 over System~1, with the filter alone accounting for $-38.1$ of the $-45.9$\cp; its separate effect on playing strength is not isolated. V3 trains a logistic gate on 18 pre-search features. Its advantage over a random gate at a 25\% search rate is $+2.5$\cp, not significant under the pre-registered test ([$-0.9$, $+5.8$], $p=0.069$) and significant under a lower-variance post hoc test ([$+0.4$, $+4.7$], Holm $p=0.02$); it recovers about 12\% of the gap to a hindsight oracle. The confidence-gating hypothesis is not supported, and the largest gains came from a component that is not a gate.
